\documentclass{ximera}

\input{../../../preamble.tex}



\definecolor{fvdnavy}{HTML}{071D43}
\definecolor{fvdcyan}{HTML}{20BFC6}
\definecolor{fvdcyanlight}{HTML}{EFFCFB}
\definecolor{fvdmagenta}{HTML}{F10D69}
\definecolor{fvdmagentalight}{HTML}{FFF1F7}
\definecolor{fvdtext}{HTML}{163044}





%=================================================
% Pedagogische omgevingen
% PDF: tcolorbox
% HTML: learning-box
%=================================================

\newenvironment{voorkennis}
{%
    \ifdefined\HCode
        \HCode{%
<section class="learning-box learning-box--prior">
<div class="learning-box__header">
<span class="learning-box__icon" aria-hidden="true"></span>
<span class="learning-box__title">Voorkennis</span>
</div>
<div class="learning-box__content">
        }%
        \begin{itemize}
    \else
       \begin{tcolorbox}[
    enhanced,
    breakable,
    colback=fvdcyanlight,
    colframe=fvdcyan,
    coltext=fvdtext,
    boxrule=0.5pt,
    leftrule=4pt,
    arc=4mm,
    outer arc=4mm,
    left=7mm,
    right=6mm,
    top=5mm,
    bottom=5mm,
    before skip=6mm,
    after skip=6mm,
    title=Voorkennis,
    coltitle=fvdcyan,
    fonttitle=\bfseries\Large,
    attach title to upper={\par\smallskip},
    boxed title style={
        empty,
        left=0mm,
        right=0mm,
        top=0mm,
        bottom=0mm
    },
    overlay={
        \draw[
            fvdcyan,
            line width=0.5pt
        ]
        ([xshift=42mm,yshift=-13mm]frame.north west)
        --
        ([xshift=-7mm,yshift=-13mm]frame.north east);
    }
]
        \begin{itemize}
    \fi
}
{%
    \end{itemize}
    \ifdefined\HCode
        \HCode{%
</div>
</section>
        }%
    \else
        \end{tcolorbox}
    \fi
}


\newenvironment{leerdoelen}
{%
    \ifdefined\HCode
        \HCode{%
<section class="learning-box learning-box--goals">
<div class="learning-box__header">
<span class="learning-box__icon" aria-hidden="true"></span>
<span class="learning-box__title">Leerdoelen</span>
</div>
<div class="learning-box__content">
        }%
        \begin{itemize}
    \else
\begin{tcolorbox}[
    enhanced,
    breakable,
    colback=fvdmagentalight,
    colframe=fvdmagenta,
    coltext=fvdtext,
    boxrule=0.5pt,
    leftrule=4pt,
    arc=4mm,
    outer arc=4mm,
    left=7mm,
    right=6mm,
    top=5mm,
    bottom=5mm,
    before skip=6mm,
    after skip=6mm,
    title=Leerdoelen,
    coltitle=fvdmagenta,
    fonttitle=\bfseries\Large,
    attach title to upper={\par\smallskip},
    boxed title style={
        empty,
        left=0mm,
        right=0mm,
        top=0mm,
        bottom=0mm
    },
    overlay={
        \draw[
            fvdmagenta,
            line width=0.5pt
        ]
        ([xshift=42mm,yshift=-13mm]frame.north west)
        --
        ([xshift=-7mm,yshift=-13mm]frame.north east);
    }
]
        \begin{itemize}
    \fi
}
{%
    \end{itemize}
    \ifdefined\HCode
        \HCode{%
</div>
</section>
        }%
    \else
        \end{tcolorbox}
    \fi
}



\addPrintStyle{../../..}

\title{Krachten en de wetten van Newton}

\author{Ingmar Herreman}

\begin{abstract}

Waarom begint een stilstaande fiets te bewegen wanneer je op de pedalen

duwt? Waarom word je in een optrekkende bus naar achteren geduwd?

Waarom trekt de aarde aan een appel, maar trekt die appel tegelijk ook

aan de aarde?

Om zulke vragen te beantwoorden hebben we het begrip kracht nodig.

Een kracht is geen eigenschap die een voorwerp zomaar bezit, maar

beschrijft een wisselwerking tussen systemen.

In dit hoofdstuk herhalen en verdiepen we de drie wetten van Newton.

We onderzoeken wat een resulterende kracht betekent, waarom een kracht

een vector is en hoe kracht en versnelling met elkaar verbonden zijn.

Bijzondere aandacht gaat naar de derde wet van Newton. Krachten ontstaan

immers steeds in wisselwerking: als voorwerp A een kracht uitoefent op

voorwerp B, dan oefent B tegelijk een even grote maar tegengesteld

gerichte kracht uit op A.

Deze ideeën vormen de brug naar de volgende hoofdstukken. Daar zullen

we twee fundamentele krachten op afstand onderzoeken: de gravitatiekracht

tussen massa's en de elektrische kracht tussen ladingen.

\end{abstract}

\begin{document}

% FVD_CHAPTER_DATA_START
% Automatisch gegenereerd uit index.tex.
% Niet handmatig aanpassen.
\fvdchapterdata{1}{Gereedschapskist en Newton - Van meten en vectoren tot krachten en beweging}{2}
% FVD_CHAPTER_DATA_END



\maketitle


% ============================================================

% VOORKENNIS EN LEERDOELEN

% ============================================================

\begin{voorkennis}

 \item Je kan scalaire en vectoriële grootheden onderscheiden.

 \item Je kan een vector beschrijven aan de hand van grootte, richting en zin.

 \item Je kan eenvoudige vectoren grafisch optellen.

 \item Je begrijpt wat een resulterende vector is.

 \item Je kent de begrippen massa, snelheid en versnelling uit de tweede graad.

 \item Je kan eenvoudige formules gebruiken en een onbekende grootheid vrijmaken.

\end{voorkennis}

\begin{leerdoelen}

 \item Je kan uitleggen dat een kracht een wisselwerking tussen twee systemen beschrijft.

 \item Je kan een kracht correct als vector voorstellen.

 \item Je kan bij eenvoudige situaties de resulterende kracht bepalen.

 \item Je kan de eerste wet van Newton uitleggen en toepassen in eenvoudige situaties.

 \item Je kan het verband tussen resulterende kracht, massa en versnelling uitleggen met de tweede wet van Newton.

 \item Je kan eenvoudige berekeningen uitvoeren met $\sum\vec F=m\vec a$.

 \item Je kan de derde wet van Newton correct formuleren.

 \item Je kan in een concrete situatie een actie-reactiepaar herkennen.

 \item Je kan uitleggen waarom actie- en reactiekracht elkaar niet opheffen.

 \item Je kan onderscheid maken tussen een krachtenpaar volgens Newton III en twee krachten die op hetzelfde voorwerp werken.

 \item Je kan het superpositiebeginsel voor krachten gebruiken.

 \item Je kan de wetten van Newton verbinden met gravitatiekracht en elektrische kracht.

\end{leerdoelen}


% ============================================================

% 1. WAAROM NEWTON?

% ============================================================

\FVDsection{Waarom verandert een beweging?}

In het vorige hoofdstuk bouwden we een aantal gereedschappen op om

fysische verschijnselen te beschrijven.

We kunnen grootheden meten, vectoren voorstellen en verschillende

vectoren combineren.

Maar daarmee is een fundamentele vraag nog niet beantwoord:

\begin{center}

\textit{Waarom verandert de beweging van een voorwerp?}

\end{center}

Een voetbal blijft liggen tot iemand ertegen trapt.

Een fiets versnelt wanneer de fietser kracht uitoefent.

Een auto vertraagt wanneer hij remt.

Een steen die wordt losgelaten versnelt naar de aarde.

De maan verandert voortdurend van bewegingsrichting terwijl ze rond

de aarde beweegt.

In al deze voorbeelden verandert de snelheidsvector.

Dat kan betekenen dat:

\begin{itemize}

 \item de grootte van de snelheid verandert;

 \item de richting van de snelheid verandert;

 \item of beide tegelijk veranderen.

\end{itemize}

Een verandering van de snelheidsvector betekent dat er een versnelling is.

\[\boxed{ \text{verandering van }\vec v \quad\Longleftrightarrow\quad \vec a\neq\vec 0 }\]

Newton legde het verband tussen zo'n bewegingsverandering en de

krachten die op een voorwerp werken.


% ============================================================

% 2. KRACHT ALS WISSELWERKING

% ============================================================

\FVDsection{Een kracht ontstaat door een wisselwerking}

We spreken in het dagelijks leven gemakkelijk over ``een kracht op

een voorwerp''.

Maar fysisch is het belangrijk om verder te vragen:

\begin{center}

\textit{Welk systeem oefent die kracht uit?}

\end{center}

Wanneer je tegen een muur duwt, oefent je hand een kracht uit op de

muur.

Wanneer de aarde een appel aantrekt, oefent de aarde een kracht uit

op de appel.

Wanneer twee geladen deeltjes elkaar afstoten, oefent elke lading een

elektrische kracht uit op de andere.


\begin{definitie}[Kracht]

Een \textbf{kracht} beschrijft een wisselwerking waarbij een systeem

een invloed uitoefent op een ander systeem.

Een kracht is een vectoriële grootheid.

\end{definitie}


We schrijven een kracht bijvoorbeeld als:

\[\vec F.\]

De SI-eenheid van kracht is de newton:

\[\boxed{ [F]=\mathrm N. }\]


\FVDsubsection{Wie oefent kracht uit op wie?}

Om verwarring te vermijden kunnen we bij een kracht expliciet

vermelden welk systeem de kracht uitoefent en op welk systeem de

kracht werkt.

We schrijven bijvoorbeeld:

\[\boxed{ \vec F_{A\rightarrow B} }\]

voor:

\begin{center}

de kracht die systeem $A$ uitoefent op systeem $B$.

\end{center}

De pijl in het onderschrift betekent dus niet de richting van de

krachtvector in de ruimte.

Hij geeft aan:

\[\boxed{ \text{wie oefent de kracht uit} \rightarrow \text{op wie werkt de kracht}. }\]


\begin{voorbeeld}[De aarde en een appel]

Een appel bevindt zich boven het aardoppervlak.

De aarde oefent een gravitatiekracht uit op de appel:

\[\vec F_{\mathrm{aarde}\rightarrow\mathrm{appel}}.\]

Maar de interactie werkt in beide richtingen.

De appel oefent eveneens een gravitatiekracht uit op de aarde:

\[\vec F_{\mathrm{appel}\rightarrow\mathrm{aarde}}.\]

Later zullen we zien dat deze twee krachten een actie-reactiepaar

vormen.

\end{voorbeeld}


\begin{oefening}[Wie oefent de kracht uit?]{\oefB\ESS}

Beschrijf bij elke situatie duidelijk welk systeem een kracht

uitoefent op welk ander systeem.

\begin{enumerate}

 \item De aarde trekt een satelliet aan.

 \item Een magneet trekt een ijzeren voorwerp aan.

 \item Een positief geladen deeltje stoot een ander positief geladen

    deeltje af.

 \item Een hand duwt tegen een deur.

 \item De vloer duwt tegen de voeten van een leerling.

\end{enumerate}

Gebruik waar mogelijk de notatie:

\[\vec F_{A\rightarrow B}.\]


\begin{oplossing}

We noteren telkens eerst het systeem dat de kracht uitoefent en daarna

het systeem waarop de kracht werkt.

\begin{enumerate}

  \item De aarde oefent een kracht uit op de satelliet:

  \[\boxed{\vec F_{\mathrm{aarde}\rightarrow\mathrm{satelliet}}}.\]

  \item De magneet oefent een kracht uit op het ijzeren voorwerp:

  \[\boxed{\vec F_{\mathrm{magneet}\rightarrow\mathrm{voorwerp}}}.\]

  \item Het eerste positief geladen deeltje oefent een afstotende

  elektrische kracht uit op het tweede:

  \[\boxed{\vec F_{1\rightarrow2}}.\]

  \item De hand oefent een kracht uit op de deur:

  \[\boxed{\vec F_{\mathrm{hand}\rightarrow\mathrm{deur}}}.\]

  \item De vloer oefent een kracht uit op de voeten van de leerling:

  \[\boxed{\vec F_{\mathrm{vloer}\rightarrow\mathrm{leerling}}}.\]

\end{enumerate}

In elk geval beschrijft de kracht een wisselwerking tussen twee systemen.

\end{oplossing}

\end{oefening}


% ============================================================

% 3. KRACHT IS VECTORIEEL

% ============================================================

\FVDsection{Een kracht is een vector}

Het is niet voldoende om te zeggen:

\begin{center}

``Er werkt een kracht van $20\,\mathrm N$''

\end{center}

We moeten ook weten in welke richting en zin de kracht werkt.

Een kracht heeft dus:

\[\boxed{ \text{grootte, richting en zin}. }\]

Daarom is kracht een vectoriële grootheid.


\begin{voorbeeld}

Twee leerlingen trekken aan een kar.

De eerste leerling oefent een kracht van:

\[50\,\mathrm N\]

naar rechts uit.

De tweede leerling oefent een kracht van:

\[30\,\mathrm N\]

naar links uit.

Omdat de krachten tegengesteld gericht zijn, is de resulterende kracht:

\[F_{\mathrm{res}} = 50-30 = 20\,\mathrm N.\]

De zin is naar rechts.

Dus:

\[\boxed{ \vec F_{\mathrm{res}} = 20\,\mathrm N \text{ naar rechts}. }\]

\end{voorbeeld}


% ============================================================

% 4. RESULTANTE

% ============================================================

\FVDsection{De resulterende kracht}

Op een voorwerp werkt meestal niet slechts één kracht.

Op een boek dat op een tafel ligt werkt bijvoorbeeld de

gravitatiekracht van de aarde, maar ook een kracht van de tafel.

Voor de beweging van het voorwerp is het gezamenlijke effect van alle

krachten belangrijk.

We noemen dit de \textbf{resulterende kracht} of \textbf{nettokracht}.


\begin{definitie}[Resulterende kracht]

De \textbf{resulterende kracht} is de vectorsom van alle krachten die

op hetzelfde gekozen systeem werken.

\[\boxed{ \vec F_{\mathrm{res}} = \sum_i\vec F_i }\]

\end{definitie}


Bij twee krachten:

\[\vec F_{\mathrm{res}} = \vec F_1+\vec F_2.\]

Bij drie krachten:

\[\vec F_{\mathrm{res}} = \vec F_1+\vec F_2+\vec F_3.\]


\begin{opmerking}

We mogen alleen krachten optellen wanneer we het effect van die

krachten op \textbf{hetzelfde systeem} onderzoeken.

Dat wordt straks bijzonder belangrijk bij de derde wet van Newton.

\end{opmerking}


\begin{oefening}[Resultante]{\oefB\ESS}

Bepaal telkens de resulterende kracht.

\begin{enumerate}

 \item $12\,\mathrm N$ naar rechts en $8\,\mathrm N$ naar rechts.

 \item $12\,\mathrm N$ naar rechts en $8\,\mathrm N$ naar links.

 \item $25\,\mathrm N$ omhoog en $25\,\mathrm N$ omlaag.

 \item $4\,\mathrm N$ naar rechts en $3\,\mathrm N$ omhoog.

\end{enumerate}

Maak telkens eerst een vectorschets.


\begin{oplossing}

Kies rechts en omhoog als positieve zin.

\begin{enumerate}

  \item Beide krachten wijzen naar rechts:

  \[F_{\mathrm{res}}=12+8=20\,\mathrm N.\]

  Dus:

  \[\boxed{\vec F_{\mathrm{res}}=20\,\mathrm N\text{ naar rechts}}.\]

  \item De krachten zijn tegengesteld:

  \[F_{\mathrm{res}}=12-8=4\,\mathrm N.\]

  Dus:

  \[\boxed{\vec F_{\mathrm{res}}=4\,\mathrm N\text{ naar rechts}}.\]

  \item De krachten zijn even groot en tegengesteld:

  \[F_{\mathrm{res}}=25-25=0.\]

  Dus:

  \[\boxed{\vec F_{\mathrm{res}}=\vec 0}.\]

  \item De krachten staan loodrecht op elkaar. Met Pythagoras:

  \[F_{\mathrm{res}}=\sqrt{4^2+3^2}=5\,\mathrm N.\]

  Voor de richting:

  \[\tan\theta=\frac{3}{4}, \qquad \theta\approx36{,}9^\circ.\]

  Dus:

  \[\boxed{F_{\mathrm{res}}=5{,}0\,\mathrm N}\]

  onder een hoek van ongeveer

  \[\boxed{36{,}9^\circ}\]

  boven de horizontale naar rechts.

\end{enumerate}

\end{oplossing}

\end{oefening}


% ============================================================

% 5. NEWTON I

% ============================================================

\FVDsection{De eerste wet van Newton: traagheid}

Stel dat een ijshockeypuck over bijna wrijvingsloos ijs glijdt.

Wat gebeurt er wanneer niemand de puck nog aanraakt?

Onze dagelijkse ervaring kan ons misleiden. Voorwerpen op aarde komen

meestal uiteindelijk tot stilstand, maar dat gebeurt omdat er

wrijvings- en weerstandskrachten aanwezig zijn.


Wanneer de resulterende kracht nul is, is er geen reden waarom de

snelheidsvector zou veranderen.


\begin{definitie}[Eerste wet van Newton]

Wanneer de resulterende kracht op een voorwerp nul is, blijft zijn

snelheidsvector constant.

\[\boxed{ \sum\vec F=\vec 0 \quad\Longrightarrow\quad \vec v=\text{constant}. }\]

\end{definitie}


Dit bevat twee belangrijke mogelijkheden.

Het voorwerp kan in rust zijn:

\[\vec v=\vec 0.\]

Maar het kan ook bewegen met constante snelheid in een rechte lijn:

\[\vec v=\text{constant}\neq\vec 0.\]


\begin{opmerking}

Een resulterende kracht van nul betekent dus \textbf{niet}

noodzakelijk dat een voorwerp stilstaat.

Het betekent dat de snelheidsvector niet verandert.

\end{opmerking}


\FVDsubsection{Traagheid}

De neiging van een voorwerp om zijn bewegingstoestand te behouden

noemen we \textbf{traagheid}.


\begin{definitie}[Traagheid]

   \begin{image}

      \includegraphics{../images/Fysica/Lambik.jpg}

      \end{image}

\textbf{Traagheid} is de eigenschap van materie waardoor een voorwerp

zich verzet tegen een verandering van zijn bewegingstoestand.

\end{definitie}

De massa is een maat voor die traagheid.

Een vrachtwagen heeft een grotere massa dan een fiets en is daarom

moeilijker in beweging te brengen, af te remmen of van richting te

veranderen.


\begin{voorbeeld}[Een optrekkende bus]

Je staat rechtop in een stilstaande bus.

Plots trekt de bus vooruit op.

Je lichaam heeft de neiging zijn oorspronkelijke bewegingstoestand

te behouden. Je voeten worden door de vloer van de bus naar voren

versneld, terwijl de rest van je lichaam niet onmiddellijk dezelfde

bewegingsverandering ondergaat.

Ten opzichte van de bus lijkt het daardoor alsof je naar achteren

beweegt.

Dit is een gevolg van traagheid.

\end{voorbeeld}


\begin{oefening}[Moet er een kracht zijn om te blijven bewegen?]{\oefU\ESS}

Een leerling zegt:

\begin{center}

"Een voorwerp kan alleen blijven bewegen zolang er een kracht in de

bewegingsrichting op werkt"

\end{center}

Beoordeel deze uitspraak met behulp van de eerste wet van Newton.

Geef ook een voorbeeld dat je redenering ondersteunt.


\begin{oplossing}

De uitspraak is fout.

Volgens de eerste wet van Newton geldt:

\[\sum\vec F=\vec0 \quad\Longrightarrow\quad \vec v=\text{constant}.\]

Een voorwerp waarop geen resulterende kracht werkt, kan dus zowel

stilstaan als met constante snelheid in een rechte lijn bewegen.

Een kracht is niet nodig om een beweging te onderhouden.

Een resulterende kracht is nodig om de snelheidsvector te

\emph{veranderen}.

Een voorbeeld is een ruimtesonde die ver van andere hemellichamen

beweegt. Wanneer de resulterende kracht verwaarloosbaar is, blijft ze

met constante snelheid in een rechte lijn bewegen zonder dat daarvoor

een kracht in de bewegingsrichting nodig is.

\end{oplossing}

\end{oefening}


% ============================================================

% 6. NEWTON II

% ============================================================

\FVDsection{Wet II - kracht verandert beweging}

De eerste wet vertelt wat gebeurt wanneer de resulterende kracht nul is.

Maar wat gebeurt er wanneer:

\[\sum\vec F\neq\vec 0?\]

Dan verandert de snelheidsvector.

Het voorwerp versnelt.


\begin{definitie}[Tweede wet van Newton]

De versnelling van een voorwerp wordt bepaald door de resulterende

kracht en de massa van het voorwerp.

\[\boxed{ \sum\vec F=m\vec a. }\]

\end{definitie}


Daaruit volgt:

\[\boxed{ \vec a = \frac{\sum\vec F}{m}. }\]


Deze formule bevat veel fysische informatie.


\FVDsubsection{Invloed van de kracht}

Bij gelijke massa geldt:

\[a\propto F_{\mathrm{res}}.\]

Een tweemaal zo grote resulterende kracht veroorzaakt een tweemaal

zo grote versnelling.


\FVDsubsection{Invloed van de massa}

Bij dezelfde resulterende kracht geldt:

\[a\propto\frac{1}{m}.\]

Een grotere massa betekent dus een kleinere versnelling bij dezelfde

resulterende kracht.

\begin{image}

   \includegraphics{../images/Fysica/newton2kracht1.png}

   \end{image}

   \begin{image}

      \includegraphics{../images/Fysica/newton2kracht2.png}

      \end{image}

\FVDsubsection{De richting van de versnelling}

Omdat $m$ een positieve scalaire grootheid is, hebben:

\[\vec a\]

en:

\[\sum\vec F\]

dezelfde richting en zin.

Dus:

\[\boxed{ \vec a\parallel\vec F_{\mathrm{res}}. }\]


\begin{opmerking}

De versnelling heeft de richting van de \textbf{resulterende kracht},

niet noodzakelijk de richting van één afzonderlijke kracht.

\end{opmerking}


\begin{voorbeeld}[Een kar versnellen]

Op een kar met massa:

\[m=4{,}0\,\mathrm{kg}\]

werkt een resulterende kracht van:

\[F_{\mathrm{res}}=12\,\mathrm N.\]

Uit:

\[F_{\mathrm{res}}=ma\]

volgt:

\[a = \frac{F_{\mathrm{res}}}{m}.\]

Dus:

\[a = \frac{12}{4{,}0} = 3{,}0\,\mathrm{m/s^2}.\]

\[\boxed{ a=3{,}0\,\mathrm{m/s^2}. }\]

\end{voorbeeld}


\begin{oefening}[Kracht, massa en versnelling]{\oefB\ESS}

\begin{enumerate}

 \item Op een voorwerp van $5{,}0\,\mathrm{kg}$ werkt een

    resulterende kracht van $20\,\mathrm N$.

    Bereken de versnelling.

 \item Een voorwerp versnelt met $3{,}0\,\mathrm{m/s^2}$ onder

    invloed van een resulterende kracht van $15\,\mathrm N$.

    Bereken de massa.

 \item Een voorwerp van $2{,}5\,\mathrm{kg}$ krijgt een versnelling

    van $4{,}0\,\mathrm{m/s^2}$.

    Bereken de resulterende kracht.

\end{enumerate}


\begin{oplossing}

We gebruiken de tweede wet van Newton:

\[F_{\mathrm{res}}=ma.\]

\begin{enumerate}

  \item

  \[a=\frac{F_{\mathrm{res}}}{m} =\frac{20}{5{,}0} =4{,}0\,\mathrm{m/s^2}.\]

  \[\boxed{a=4{,}0\,\mathrm{m/s^2}}.\]

  \item

  \[m=\frac{F_{\mathrm{res}}}{a} =\frac{15}{3{,}0} =5{,}0\,\mathrm{kg}.\]

  \[\boxed{m=5{,}0\,\mathrm{kg}}.\]

  \item

  \[F_{\mathrm{res}}=ma =(2{,}5)(4{,}0) =10\,\mathrm N.\]

  \[\boxed{F_{\mathrm{res}}=10\,\mathrm N}.\]

\end{enumerate}

\end{oplossing}

\end{oefening}


\begin{oefening}[Dezelfde kracht, andere massa]{\oefU\ESS}

Twee karretjes worden afzonderlijk voortgeduwd door dezelfde constante

resulterende kracht.

Karretje A heeft massa:

\[m_A=2{,}0\,\mathrm{kg}.\]

Karretje B heeft massa:

\[m_B=6{,}0\,\mathrm{kg}.\]

Zonder de kracht numeriek te kennen:

\begin{enumerate}

 \item vergelijk $a_A$ en $a_B$;

 \item bepaal de verhouding

 \[\frac{a_A}{a_B};\]

 \item verklaar je antwoord fysisch met het begrip traagheid.

\end{enumerate}


\begin{oplossing}

Voor beide karretjes is de resulterende kracht dezelfde.

Volgens Newton II:

\[a=\frac{F_{\mathrm{res}}}{m}.\]

Dus:

\[a_A=\frac{F}{2{,}0} \qquad\text{en}\qquad a_B=\frac{F}{6{,}0}.\]

\begin{enumerate}

  \item Omdat karretje A de kleinste massa heeft:

  \[\boxed{a_A>a_B}.\]

  \item

  \[\frac{a_A}{a_B} = \frac{F/m_A}{F/m_B} = \frac{m_B}{m_A} = \frac{6{,}0}{2{,}0} =3.\]

  Dus:

  \[\boxed{\frac{a_A}{a_B}=3}.\]

  Karretje A versnelt driemaal zo sterk als karretje B.

  \item Karretje B heeft een driemaal zo grote massa en dus een grotere

  traagheid. Bij dezelfde resulterende kracht verandert zijn

  bewegingstoestand daarom minder sterk.

\end{enumerate}

\end{oplossing}

\end{oefening}


% ============================================================

% 7. NEWTON III

% ============================================================

\FVDsection{Wet III - krachten komen in paren}

Tot nu toe bekeken we vooral welke krachten op één gekozen voorwerp

werken.

Maar een kracht ontstaat altijd door een interactie tussen twee

systemen.

Wanneer systeem $A$ een kracht uitoefent op systeem $B$, oefent

systeem $B$ tegelijk een kracht uit op systeem $A$.


\begin{definitie}[Derde wet van Newton]

Wanneer systeem $A$ een kracht uitoefent op systeem $B$, oefent

systeem $B$ tegelijkertijd een even grote maar tegengesteld gerichte

kracht uit op systeem $A$.

\[\boxed{ \vec F_{A\rightarrow B} = -\vec F_{B\rightarrow A}. }\]

\end{definitie}


De twee krachten:

\begin{itemize}

 \item ontstaan door dezelfde interactie;

 \item zijn even groot;

 \item hebben dezelfde richting;

 \item hebben tegengestelde zin;

 \item treden gelijktijdig op;

 \item werken op verschillende systemen.

\end{itemize}


\begin{opmerking}

Het laatste punt is cruciaal:

\[\boxed{ \text{actie en reactie werken op verschillende voorwerpen.} }\]

Daarom heffen ze elkaar niet op wanneer we de beweging van één

afzonderlijk voorwerp onderzoeken.

\end{opmerking}

\begin{image}

   \includegraphics{../images/Fysica/newton3.png}

   \end{image}

% ============================================================

% 8. ACTIE EN REACTIE

% ============================================================

\FVDsubsection{Actie en reactie herkennen}

De woorden ``actie'' en ``reactie'' kunnen misleidend zijn.

Ze betekenen niet dat eerst de ene kracht ontstaat en pas daarna de

andere.

Beide krachten bestaan tegelijkertijd.

Er is ook geen fysisch onderscheid waarbij één kracht belangrijker

of ``de echte kracht'' zou zijn.

Beide maken deel uit van dezelfde interactie.


\begin{voorbeeld}[Hand en muur]

Je duwt met je hand tegen een muur.

Je hand oefent een kracht uit op de muur:

\[\vec F_{\mathrm{hand}\rightarrow\mathrm{muur}}.\]

De muur oefent tegelijkertijd een kracht uit op je hand:

\[\vec F_{\mathrm{muur}\rightarrow\mathrm{hand}}.\]

Volgens Newton III:

\[\boxed{ \vec F_{\mathrm{hand}\rightarrow\mathrm{muur}} = - \vec F_{\mathrm{muur}\rightarrow\mathrm{hand}}. }\]

Daarom voel je de muur tegen je hand duwen.

\end{voorbeeld}


\begin{voorbeeld}[Twee skaters]

Twee skaters staan stil tegenover elkaar op vrijwel wrijvingsloos ijs.

Skater A duwt tegen skater B.

Dan oefent A een kracht uit op B:

\[\vec F_{A\rightarrow B}.\]

Tegelijk oefent B een kracht uit op A:

\[\vec F_{B\rightarrow A}.\]

De krachten zijn even groot:

\[F_{A\rightarrow B} = F_{B\rightarrow A}.\]

Toch hoeven hun versnellingen niet even groot te zijn.

Volgens Newton II:

\[a_A = \frac{F}{m_A}\]

en:

\[a_B = \frac{F}{m_B}.\]

Als:

\[m_A<m_B,\]

dan:

\[a_A>a_B.\]

De lichtere skater krijgt dus de grootste versnelling, hoewel beide

krachten even groot zijn.

\end{voorbeeld}


\begin{oefening}[Twee skaters]{\oefU\ESS}

Skater A heeft massa:

\[m_A=50\,\mathrm{kg}\]

en skater B:

\[m_B=75\,\mathrm{kg}.\]

Tijdens het afduwen oefent A een kracht van:

\[120\,\mathrm N\]

uit op B.

\begin{question}

Hoe groot is de kracht die B op A uitoefent?

\end{question}

\begin{question}

Bereken de versnelling van A.

\end{question}

\begin{question}

Bereken de versnelling van B.

\end{question}

\begin{question}

Leg uit waarom de versnellingen verschillend zijn terwijl de krachten

even groot zijn.

\end{question}


\begin{oplossing}

Volgens Newton III zijn de krachten die beide skaters op elkaar

uitoefenen even groot en tegengesteld gericht.

\begin{enumerate}

  \item B oefent een kracht van

  \[\boxed{120\,\mathrm N}\]

  uit op A, in de tegengestelde zin van de kracht van A op B.

  \item Voor skater A:

  \[a_A=\frac{F}{m_A} =\frac{120}{50} =2{,}4\,\mathrm{m/s^2}.\]

  Dus:

  \[\boxed{a_A=2{,}4\,\mathrm{m/s^2}}.\]

  \item Voor skater B:

  \[a_B=\frac{F}{m_B} =\frac{120}{75} =1{,}6\,\mathrm{m/s^2}.\]

  Dus:

  \[\boxed{a_B=1{,}6\,\mathrm{m/s^2}}.\]

  \item De krachten zijn volgens Newton III even groot, maar de massa's

  zijn verschillend. Volgens

  \[a=\frac{F}{m}\]

  krijgt bij dezelfde kracht de kleinste massa de grootste versnelling.

  Daarom:

  \[\boxed{a_A>a_B}.\]

\end{enumerate}

\end{oplossing}

\end{oefening}


% ============================================================

% 9. ACTIE-REACTIE VERSUS EVENWICHT

% ============================================================

\FVDsubsection{Een belangrijke valkuil: krachtpaar of krachtenevenwicht?}

Dit is een van de belangrijkste onderscheidingen in de mechanica.

Beschouw een boek dat stil op een tafel ligt.

Op het boek werkt de gravitatiekracht van de aarde:

\[\vec F_{\mathrm{aarde}\rightarrow\mathrm{boek}}.\]

Daarnaast oefent de tafel een opwaartse kracht uit op het boek:

\[\vec F_{\mathrm{tafel}\rightarrow\mathrm{boek}}.\]

Wanneer het boek in rust blijft, kunnen deze twee krachten elkaar

compenseren:

\[\vec F_{\mathrm{aarde}\rightarrow\mathrm{boek}} + \vec F_{\mathrm{tafel}\rightarrow\mathrm{boek}} = \vec 0.\]

Maar deze twee krachten vormen \textbf{geen} actie-reactiepaar.

Waarom niet?

Omdat beide krachten op hetzelfde systeem werken:

\[\boxed{\text{het boek}.}\]


De reactiekracht bij:

\[\vec F_{\mathrm{aarde}\rightarrow\mathrm{boek}}\]

is:

\[\vec F_{\mathrm{boek}\rightarrow\mathrm{aarde}}.\]

En de reactiekracht bij:

\[\vec F_{\mathrm{tafel}\rightarrow\mathrm{boek}}\]

is:

\[\vec F_{\mathrm{boek}\rightarrow\mathrm{tafel}}.\]


\begin{opmerking}

Gebruik deze test:

\[\boxed{ \vec F_{A\rightarrow B} \quad\leftrightarrow\quad \vec F_{B\rightarrow A} }\]

Bij een actie-reactiepaar worden de namen van de twee systemen dus

omgewisseld.

\end{opmerking}


\begin{oefening}[Actie-reactie of niet?]{\oefB\ESS}

Bepaal telkens of de twee beschreven krachten een actie-reactiepaar

vormen. Verklaar.

\begin{enumerate}

 \item De aarde trekt aan een appel en de appel trekt aan de aarde.

 \item De aarde trekt een boek naar beneden en de tafel duwt het boek

    omhoog.

 \item Een voet duwt tegen een voetbal en de voetbal duwt tegen de voet.

 \item Een auto trekt een aanhangwagen vooruit en de aanhangwagen trekt

    de auto achteruit.

 \item De aarde trekt een satelliet aan en de zon trekt dezelfde

    satelliet aan.

\end{enumerate}


\begin{oplossing}

\begin{enumerate}

  \item \textbf{Ja.}

  \[\vec F_{\mathrm{aarde}\rightarrow\mathrm{appel}} = -\vec F_{\mathrm{appel}\rightarrow\mathrm{aarde}}.\]

  De twee krachten horen bij dezelfde interactie en werken op

  verschillende systemen.

  \item \textbf{Nee.}

  Beide genoemde krachten werken op het boek. Ze kunnen elkaar

  compenseren, maar vormen geen actie-reactiepaar.

  \item \textbf{Ja.}

  \[\vec F_{\mathrm{voet}\rightarrow\mathrm{bal}} = -\vec F_{\mathrm{bal}\rightarrow\mathrm{voet}}.\]

  \item \textbf{Ja.}

  \[\vec F_{\mathrm{auto}\rightarrow\mathrm{aanhangwagen}} = -\vec F_{\mathrm{aanhangwagen}\rightarrow\mathrm{auto}}.\]

  \item \textbf{Nee.}

  Beide krachten werken op dezelfde satelliet en zijn afkomstig van

  twee verschillende interacties.

  De reactiekracht op de kracht van de aarde op de satelliet is de

  kracht van de satelliet op de aarde. De reactiekracht op de kracht

  van de zon op de satelliet is de kracht van de satelliet op de zon.

\end{enumerate}

\end{oplossing}

\end{oefening}


% ============================================================

% 10. DE RAKET

% ============================================================

\FVDsection{Hoe kan een raket zich afzetten in de ruimte?}

Een hardnekkige misconceptie is dat een raket zich in de ruimte niet

zou kunnen voortbewegen omdat er ``niets is om zich tegen af te zetten''.

Maar een raket hoeft zich niet af te zetten tegen de lucht.

De interactie vindt plaats tussen de raket en de uitgestoten gassen.

De raket oefent een kracht uit op de uitlaatgassen:

\[\vec F_{\mathrm{raket}\rightarrow\mathrm{gas}}.\]

De gassen oefenen tegelijkertijd een tegengestelde kracht uit op de

raket:

\[\vec F_{\mathrm{gas}\rightarrow\mathrm{raket}}.\]

Volgens Newton III:

\[\boxed{ \vec F_{\mathrm{raket}\rightarrow\mathrm{gas}} = - \vec F_{\mathrm{gas}\rightarrow\mathrm{raket}}. }\]

De raket kan daardoor ook in het luchtledige versnellen.


\begin{oefening}[Raket in het luchtledige]{\oefU\ESS}

Een leerling beweert:

\begin{center} ``Een raket kan in de ruimte niet versnellen, want er is geen lucht

waartegen hij zich kan afzetten.''

\end{center}

Leg met behulp van Newton III uit waarom deze redenering fout is.

Vermeld expliciet de twee systemen en de twee krachten van het

actie-reactiepaar.


\begin{oplossing}

De redenering is fout. Een raket hoeft zich niet tegen de lucht af te

zetten.

De twee relevante systemen zijn:

\begin{itemize}

  \item de raket;

  \item de uitgestoten gassen.

\end{itemize}

De raket oefent een kracht uit op de gassen:

\[\vec F_{\mathrm{raket}\rightarrow\mathrm{gas}}.\]

Tegelijk oefenen de gassen een kracht uit op de raket:

\[\vec F_{\mathrm{gas}\rightarrow\mathrm{raket}}.\]

Volgens Newton III:

\[\boxed{ \vec F_{\mathrm{raket}\rightarrow\mathrm{gas}} = -\vec F_{\mathrm{gas}\rightarrow\mathrm{raket}} }.\]

De raket versnelt door de kracht die de uitgestoten gassen op de raket

uitoefenen. Daarvoor is geen atmosfeer nodig.

\end{oplossing}

\end{oefening}


% ============================================================

% 11. SUPERPOSITIE

% ============================================================

\FVDsection{Meerdere krachten tegelijk: superpositie}

Een voorwerp kan tegelijkertijd met meerdere andere systemen

interageren.

Een satelliet kan bijvoorbeeld worden aangetrokken door de aarde,

maar ook door de maan en de zon.

Een geladen deeltje kan tegelijkertijd elektrische krachten

ondervinden van meerdere andere ladingen.

Het gezamenlijke effect vinden we door de krachtvectoren op te tellen.


\begin{definitie}[Superpositiebeginsel voor krachten]

Wanneer meerdere krachten op hetzelfde systeem werken, is de

resulterende kracht de vectorsom van de afzonderlijke krachten.

\[\boxed{ \vec F_{\mathrm{res}} = \sum_i\vec F_i. }\]

\end{definitie}


Voor drie krachten:

\[\boxed{ \vec F_{\mathrm{res}} = \vec F_1+\vec F_2+\vec F_3. }\]


\begin{opmerking}

Superpositie is straks essentieel.

Bij gravitatie zullen verschillende massa's elk een gravitatiekracht

uitoefenen.

Bij elektriciteit zullen verschillende ladingen elk een elektrische

kracht uitoefenen.

De totale kracht vinden we telkens door de afzonderlijke

\textbf{krachtvectoren} op te tellen.

\end{opmerking}


\begin{voorbeeld}[Twee tegengestelde krachten]

Op een deeltje werken twee krachten langs dezelfde rechte:

\[F_1=8{,}0\,\mathrm N\]

naar rechts en:

\[F_2=3{,}0\,\mathrm N\]

naar links.

Kiezen we rechts als positieve zin, dan:

\[F_{\mathrm{res}} = 8{,}0-3{,}0 = 5{,}0\,\mathrm N.\]

Dus:

\[\boxed{ \vec F_{\mathrm{res}} = 5{,}0\,\mathrm N \text{ naar rechts}. }\]

\end{voorbeeld}


\begin{oefening}[Twee loodrechte krachten]{\oefU\ESS}

Op een klein voorwerp werken twee krachten:

\[F_1=6{,}0\,\mathrm N\]

naar rechts en:

\[F_2=8{,}0\,\mathrm N\]

omhoog.

\begin{question}

Maak een vectorschets.

\end{question}

\begin{question}

Bereken de grootte van de resulterende kracht.

\end{question}

\begin{question}

Bepaal de richting van de resulterende kracht.

\end{question}

\begin{question}

Het voorwerp heeft een massa van $2{,}0\,\mathrm{kg}$.

Bereken de grootte van zijn versnelling.

\end{question}

\begin{question}

Welke richting heeft de versnelling?

\end{question}


\begin{oplossing}

De horizontale en verticale kracht staan loodrecht op elkaar.

\begin{enumerate}

  \item De vectorschets bevat een krachtvector van

  $6{,}0\,\mathrm N$ naar rechts en vanuit hetzelfde aangrijpingspunt

  een krachtvector van $8{,}0\,\mathrm N$ omhoog. De resultante wijst

  schuin naar rechtsboven.

  \item Met Pythagoras:

  \[F_{\mathrm{res}} = \sqrt{6{,}0^2+8{,}0^2} = 10{,}0\,\mathrm N.\]

  Dus:

  \[\boxed{F_{\mathrm{res}}=10{,}0\,\mathrm N}.\]

  \item Noem $\theta$ de hoek met de horizontale naar rechts.

  \[\tan\theta=\frac{8{,}0}{6{,}0},\]

  zodat

  \[\theta=\arctan\left(\frac{8}{6}\right) \approx53{,}1^\circ.\]

  De resultante wijst dus

  \[\boxed{53{,}1^\circ}\]

  boven de horizontale naar rechts.

  \item Volgens Newton II:

  \[a=\frac{F_{\mathrm{res}}}{m} =\frac{10{,}0}{2{,}0} =5{,}0\,\mathrm{m/s^2}.\]

  Dus:

  \[\boxed{a=5{,}0\,\mathrm{m/s^2}}.\]

  \item Omdat

  \[\vec a=\frac{\vec F_{\mathrm{res}}}{m},\]

  heeft de versnelling dezelfde richting en zin als de resulterende

  kracht: $53{,}1^\circ$ boven de horizontale naar rechts.

\end{enumerate}

\end{oplossing}

\end{oefening}


% ============================================================

% 12. NEWTON I EN II SAMEN

% ============================================================

\FVDsection{Newton I en Newton II vertellen hetzelfde verhaal}

De eerste en tweede wet van Newton worden vaak als volledig aparte

wetten geleerd.

Toch sluiten ze sterk bij elkaar aan.

Volgens Newton II:

\[\sum\vec F=m\vec a.\]

Wanneer:

\[\sum\vec F=\vec 0,\]

dan volgt:

\[m\vec a=\vec 0.\]

Omdat:

\[m\neq0,\]

moet:

\[\vec a=\vec 0.\]

Dus verandert de snelheidsvector niet:

\[\boxed{ \vec v=\text{constant}. }\]

Dat is precies de bewegingstoestand die in Newton I wordt beschreven.

Newton I benadrukt daarnaast het fundamentele idee van traagheid en

de bijzondere betekenis van inertiaalstelsels, waarop we later nog

kunnen terugkomen.


% ============================================================

% 13. VAN KRACHT NAAR BEWEGING

% ============================================================

\FVDsection{Van interactie naar bewegingsverandering}

We kunnen de ideeën uit dit hoofdstuk nu samenbrengen.

Wanneer twee systemen met elkaar interageren, ontstaan krachten.

Volgens Newton III:

\[\boxed{ \vec F_{A\rightarrow B} = -\vec F_{B\rightarrow A}. }\]

Op één systeem kunnen verschillende krachten werken.

Door superpositie:

\[\boxed{ \vec F_{\mathrm{res}} = \sum_i\vec F_i. }\]

De resulterende kracht bepaalt vervolgens de versnelling:

\[\boxed{ \sum\vec F=m\vec a. }\]

We krijgen dus de keten:

\[\boxed{ \text{interactie} \longrightarrow \text{krachten} \longrightarrow \text{resulterende kracht} \longrightarrow \text{versnelling} \longrightarrow \text{verandering van beweging}. }\]


% ============================================================

% 14. EEN DIEPERE VRAAG

% ============================================================

\FVDsection{Even grote krachten, maar niet hetzelfde effect}

Newton III zegt:

\[F_{A\rightarrow B} = F_{B\rightarrow A}.\]

Daaruit zou je ten onrechte kunnen besluiten dat beide voorwerpen ook

even sterk versnellen.

Dat hoeft niet.

Voor systeem A:

\[a_A = \frac{F_{B\rightarrow A}}{m_A}.\]

Voor systeem B:

\[a_B = \frac{F_{A\rightarrow B}}{m_B}.\]

Omdat de groottes van de krachten gelijk zijn:

\[F_{A\rightarrow B} = F_{B\rightarrow A} = F,\]

geldt:

\[a_A=\frac{F}{m_A}\]

en:

\[a_B=\frac{F}{m_B}.\]

Daaruit volgt:

\[\boxed{ \frac{a_A}{a_B} = \frac{m_B}{m_A}. }\]

De versnellingen zijn dus omgekeerd evenredig met de massa's.


\begin{voorbeeld}[Appel en aarde]

De aarde trekt aan een vallende appel.

Maar volgens Newton III trekt de appel met precies dezelfde kracht

aan de aarde.

Waarom zien we de aarde dan niet merkbaar naar de appel versnellen?

Omdat:

\[m_{\mathrm{aarde}}\gg m_{\mathrm{appel}}.\]

Voor dezelfde kracht geldt:

\[a=\frac{F}{m}.\]

De versnelling van de aarde is daardoor extreem klein.

\begin{center}

\textbf{De aarde trekt dus niet harder aan de appel dan de appel aan de aarde.}

\end{center}

Het verschil zit in de massa's en daardoor in de versnellingen.

\end{voorbeeld}


\begin{oefening}[Aarde en astronaut]{\oefU\ESS}

Een astronaut bevindt zich in de ruimte en wordt door de aarde

aangetrokken.

Een leerling zegt:

\begin{center} ``De aarde oefent een veel grotere kracht uit op de astronaut dan de

astronaut op de aarde, want de aarde heeft veel meer massa.''

\end{center}

\begin{question}

Is deze uitspraak juist?

\end{question}

\begin{question}

Gebruik Newton III in je verklaring.

\end{question}

\begin{question}

Waarom is de versnelling van de astronaut toch veel groter dan die van

de aarde?

\end{question}


\begin{oplossing}

\begin{enumerate}

  \item Nee. De uitspraak is fout.

  \item Volgens Newton III:

  \[\boxed{ \vec F_{\mathrm{aarde}\rightarrow\mathrm{astronaut}} = -\vec F_{\mathrm{astronaut}\rightarrow\mathrm{aarde}} }.\]

  De aarde en de astronaut oefenen dus krachten met dezelfde grootte

  en tegengestelde zin op elkaar uit.

  \item Voor beide systemen geldt Newton II:

  \[a=\frac{F}{m}.\]

  Hoewel de krachtgroottes gelijk zijn, is

  \[m_{\mathrm{aarde}}\gg m_{\mathrm{astronaut}}.\]

  Daarom:

  \[a_{\mathrm{aarde}} = \frac{F}{m_{\mathrm{aarde}}} \ll \frac{F}{m_{\mathrm{astronaut}}} = a_{\mathrm{astronaut}}.\]

  De astronaut krijgt dus een veel grotere versnelling door zijn veel

  kleinere massa.

\end{enumerate}

\end{oplossing}

\end{oefening}


\begin{oefening}[Waar of niet waar?]{\oefB\ESS}

Beoordeel elke uitspraak. Verbeter een foute uitspraak.

\begin{enumerate}

 \item Als een voorwerp beweegt, moet er een resulterende kracht in

    de bewegingsrichting werken.

 \item Als de resulterende kracht nul is, moet het voorwerp stilstaan.

 \item Als de resulterende kracht nul is, is de versnelling nul.

 \item De versnelling heeft dezelfde richting als de resulterende

    kracht.

 \item Een grotere massa betekent bij dezelfde resulterende kracht

    een grotere versnelling.

 \item Actie en reactie werken op hetzelfde voorwerp.

 \item Actie en reactie zijn even groot.

 \item Actie en reactie ontstaan tegelijkertijd.

\end{enumerate}


\begin{oplossing}

\begin{enumerate}

  \item \textbf{Niet waar.}

  Een voorwerp kan met constante snelheid bewegen terwijl

  \[\sum\vec F=\vec0.\]

  \item \textbf{Niet waar.}

  Een resulterende kracht van nul betekent dat de snelheid constant is.

  Die constante snelheid kan ook verschillend van nul zijn.

  \item \textbf{Waar.}

  Uit

  \[\sum\vec F=m\vec a\]

  volgt bij $\sum\vec F=\vec0$ dat

  \[\vec a=\vec0.\]

  \item \textbf{Waar.}

  Omdat de massa positief is, hebben $\vec a$ en

  $\vec F_{\mathrm{res}}$ dezelfde richting en zin.

  \item \textbf{Niet waar.}

  Bij dezelfde kracht geldt

  \[a=\frac{F}{m}.\]

  Een grotere massa geeft dus een \emph{kleinere} versnelling.

  \item \textbf{Niet waar.}

  Actie en reactie werken altijd op verschillende systemen.

  \item \textbf{Waar.}

  Volgens Newton III zijn de krachtgroottes gelijk.

  \item \textbf{Waar.}

  De twee krachten van een actie-reactiepaar bestaan gelijktijdig.

\end{enumerate}

\end{oplossing}

\end{oefening}


\begin{oefening}[Welke wet van Newton?]{\oefB\ESS}

Geef bij elke situatie aan welke wet van Newton het meest rechtstreeks

wordt geïllustreerd. Verklaar kort.

\begin{enumerate}

 \item Een passagier beweegt bij plots remmen naar voren ten opzichte

    van de auto.

 \item Een leeg winkelkarretje krijgt bij dezelfde duwkracht een

    grotere versnelling dan een vol winkelkarretje.

 \item Een zwemmer duwt water naar achteren en beweegt daardoor naar

    voren.

 \item Een ruimtesonde blijft met constante snelheid bewegen wanneer

    de resulterende kracht erop nul is.

 \item Een geweer ondervindt bij het afvuren een terugslag.

\end{enumerate}


\begin{oplossing}

\begin{enumerate}

  \item \textbf{Newton I.}

  Bij het remmen heeft de passagier door traagheid de neiging zijn

  oorspronkelijke bewegingstoestand te behouden.

  \item \textbf{Newton II.}

  Uit

  \[a=\frac{F_{\mathrm{res}}}{m}\]

  volgt dat bij dezelfde kracht de kleinste massa de grootste

  versnelling krijgt.

  \item \textbf{Newton III.}

  De zwemmer duwt het water naar achteren en het water oefent

  tegelijkertijd een tegengestelde kracht uit op de zwemmer.

  \item \textbf{Newton I.}

  Wanneer

  \[\sum\vec F=\vec0,\]

  blijft de snelheid constant.

  \item \textbf{Newton III.}

  Het geweer oefent een kracht uit op de kogel en de kogel oefent

  tegelijkertijd een tegengestelde kracht uit op het geweer.

\end{enumerate}

\end{oplossing}

\end{oefening}


\begin{oefening}[De parachutist]{\oefU\ESS}

Een parachutist valt met constante snelheid naar beneden.

Op de parachutist werken onder andere:

\begin{itemize}

 \item een neerwaartse gravitatiekracht;

 \item een opwaartse luchtweerstandskracht.

\end{itemize}

\begin{question}

Wat kan je besluiten over de resulterende kracht?

\end{question}

\begin{question}

Is de versnelling nul? Verklaar.

\end{question}

\begin{question}

Werken er geen krachten omdat de snelheid constant is?

\end{question}

\begin{question}

Vormen de gravitatiekracht en de luchtweerstandskracht een

actie-reactiepaar? Verklaar.

\end{question}


\begin{oplossing}

\begin{enumerate}

  \item De parachutist beweegt met constante snelheid, dus:

  \[\vec a=\vec0.\]

  Volgens Newton II:

  \[\sum\vec F=m\vec a=\vec0.\]

  Dus:

  \[\boxed{\vec F_{\mathrm{res}}=\vec0}.\]

  \item Ja. De versnelling is nul omdat de snelheidsvector constant is:

  \[\boxed{\vec a=\vec0}.\]

  \item Nee. Er werken wel degelijk krachten. De neerwaartse

  gravitatiekracht en de opwaartse luchtweerstand zijn even groot en

  tegengesteld gericht. Hun vectorsom is nul:

  \[\boxed{ \vec F_{\mathrm g}+\vec F_{\mathrm{lucht}}=\vec0 }.\]

  \item Nee. Beide genoemde krachten werken op dezelfde parachutist en

  kunnen daarom geen actie-reactiepaar vormen.

  De reactiekracht op de gravitatiekracht van de aarde op de parachutist

  is de gravitatiekracht van de parachutist op de aarde. De

  reactiekracht op de kracht van de lucht op de parachutist is de kracht

  van de parachutist op de lucht.

\end{enumerate}

\end{oplossing}

\end{oefening}


\begin{opmerking}

Een parachutist die met constante snelheid daalt noemen we hier niet ``in vrije val''. In het ideale vrije-valmodel verwaarlozen we de

luchtweerstand en werkt alleen de gravitatiekracht.

\end{opmerking}


% ============================================================

% 18. ANALYSE-OEFENINGEN

% ============================================================

\begin{oefening}[Een botsing tussen een vrachtwagen en een auto]{\oefU\ESS}

Een zware vrachtwagen botst frontaal met een veel lichtere personenauto.

Een leerling zegt:

\begin{center} ``De vrachtwagen oefent tijdens de botsing een grotere kracht uit op

de auto, want de vrachtwagen is veel zwaarder.''

\end{center}

Analyseer deze redenering.

Bespreek afzonderlijk:

\begin{enumerate}

 \item de kracht van de vrachtwagen op de auto;

 \item de kracht van de auto op de vrachtwagen;

 \item de versnelling van de vrachtwagen;

 \item de versnelling van de auto.

\end{enumerate}

Gebruik Newton II en Newton III in je verklaring.


\begin{oplossing}

De uitspraak van de leerling is fout.

Volgens Newton III geldt tijdens de botsing:

\[\boxed{ \vec F_{\mathrm{vrachtwagen}\rightarrow\mathrm{auto}} = -\vec F_{\mathrm{auto}\rightarrow\mathrm{vrachtwagen}} }.\]

\begin{enumerate}

  \item De vrachtwagen oefent een kracht uit op de auto.

  \item De auto oefent tegelijkertijd een even grote, maar tegengesteld

  gerichte kracht uit op de vrachtwagen. De vrachtwagen oefent dus

  \emph{geen grotere kracht} uit omdat hij zwaarder is.

  \item Voor de vrachtwagen geldt:

  \[a_{\mathrm v}=\frac{F}{m_{\mathrm v}}.\]

  \item Voor de auto geldt:

  \[a_{\mathrm a}=\frac{F}{m_{\mathrm a}}.\]

  Omdat

  \[m_{\mathrm v}>m_{\mathrm a},\]

  volgt bij dezelfde krachtgrootte:

  \[|a_{\mathrm v}|<|a_{\mathrm a}|.\]

\end{enumerate}

De lichtere auto ondervindt dus een grotere versnelling in grootte,

niet een grotere interactiekracht.

Newton III bepaalt de relatie tussen de twee krachten; Newton II

verklaart waarom dezelfde krachtgroottes verschillende versnellingen

kunnen veroorzaken.

\end{oplossing}

\end{oefening}


\begin{oefening}[Touwtrekken]{\oefU}

Twee leerlingen trekken aan de uiteinden van hetzelfde touw.

Leerling A beweert:

\begin{center} ``Ik trek harder aan het touw dan jij aan mij trekt, want ik win.''

\end{center}

Onderzoek zorgvuldig welke krachten een actie-reactiepaar vormen.

Leg vervolgens uit waarom het feit dat één leerling beweegt niet

betekent dat Newton III geschonden wordt.


\begin{oplossing}

De uitspraak is niet correct als daarmee een actie-reactiepaar wordt

bedoeld.

Wanneer leerling A aan het touw trekt, oefent A een kracht uit op het

touw:

\[\vec F_{A\rightarrow\mathrm{touw}}.\]

Het bijbehorende actie-reactiepaar is:

\[\vec F_{\mathrm{touw}\rightarrow A}.\]

Volgens Newton III:

\[\vec F_{A\rightarrow\mathrm{touw}} = -\vec F_{\mathrm{touw}\rightarrow A}.\]

Aan de andere kant vormen

\[\vec F_{B\rightarrow\mathrm{touw}} \quad\text{en}\quad \vec F_{\mathrm{touw}\rightarrow B}\]

een ander actie-reactiepaar.

Dat één leerling uiteindelijk beweegt, betekent niet dat Newton III

geschonden wordt. Voor de beweging van een leerling moeten we alle

krachten optellen die \emph{op die leerling} werken. Naast de kracht

van het touw spelen bijvoorbeeld de krachten tussen de voeten en de

grond een rol.

Newton III vergelijkt krachten die op \emph{verschillende} systemen

werken. Newton II bepaalt uit de resulterende kracht op één systeem

hoe dat systeem versnelt.

\end{oplossing}

\end{oefening}


\begin{oefening}[Drie krachten op een ruimtesonde]{\oefU\ESS}

Op een ruimtesonde werken op een bepaald ogenblik drie krachten.

\[\vec F_1: \quad 6{,}0\,\mathrm N \text{ naar rechts},\]

\[\vec F_2: \quad 2{,}0\,\mathrm N \text{ naar links},\]

\[\vec F_3: \quad 3{,}0\,\mathrm N \text{ omhoog}.\]

De massa van de sonde is:

\[m=2{,}0\,\mathrm{kg}.\]

\begin{question}

Maak een vectorschets.

\end{question}

\begin{question}

Bepaal de horizontale en verticale component van de resulterende kracht.

\end{question}

\begin{question}

Bereken de grootte van de resulterende kracht.

\end{question}

\begin{question}

Bepaal de richting van de resulterende kracht.

\end{question}

\begin{question}

Bereken de grootte van de versnelling.

\end{question}

\begin{question}

Welke richting heeft de versnelling?

\end{question}


\begin{oplossing}

Kies rechts als positieve $x$-richting en omhoog als positieve

$y$-richting.

\begin{enumerate}

  \item De vectorschets bevat $6{,}0\,\mathrm N$ naar rechts,

  $2{,}0\,\mathrm N$ naar links en $3{,}0\,\mathrm N$ omhoog.

  De resultante wijst naar rechtsboven.

  \item Horizontaal:

  \[F_{\mathrm{res},x}=6{,}0-2{,}0=4{,}0\,\mathrm N.\]

  Verticaal:

  \[F_{\mathrm{res},y}=3{,}0\,\mathrm N.\]

  Dus:

  \[\boxed{ F_{\mathrm{res},x}=4{,}0\,\mathrm N, \qquad F_{\mathrm{res},y}=3{,}0\,\mathrm N }.\]

  \item Met Pythagoras:

  \[F_{\mathrm{res}} = \sqrt{4{,}0^2+3{,}0^2} = 5{,}0\,\mathrm N.\]

  \[\boxed{F_{\mathrm{res}}=5{,}0\,\mathrm N}.\]

  \item Voor de hoek $\theta$ boven de horizontale:

  \[\tan\theta=\frac{3{,}0}{4{,}0},\]

  dus

  \[\theta\approx36{,}9^\circ.\]

  \item Volgens Newton II:

  \[a=\frac{F_{\mathrm{res}}}{m} =\frac{5{,}0}{2{,}0} =2{,}5\,\mathrm{m/s^2}.\]

  \[\boxed{a=2{,}5\,\mathrm{m/s^2}}.\]

  \item De versnelling heeft dezelfde richting en zin als de

  resulterende kracht:

  \[\boxed{36{,}9^\circ}\]

  boven de horizontale naar rechts.

\end{enumerate}

\end{oplossing}

\end{oefening}


% ============================================================

% 19. GEVORDERDE OEFENINGEN

% ============================================================

\begin{oefening}[Drie massa's op één lijn]{\oefG}

Drie massa's $A$, $B$ en $C$ liggen op één rechte.

Massa $A$ bevindt zich links van $B$ en massa $C$ rechts van $B$.

Door gravitatie oefent $A$ een kracht uit op $B$:

\[\vec F_{A\rightarrow B}.\]

Ook $C$ oefent een gravitatiekracht uit op $B$:

\[\vec F_{C\rightarrow B}.\]

Veronderstel dat:

\[F_{A\rightarrow B} = F_{C\rightarrow B}.\]

\begin{question}

Bepaal de resulterende gravitatiekracht op $B$.

\end{question}

\begin{question}

Wat kan je op basis van Newton II zeggen over de versnelling van $B$

op dat ogenblik?

\end{question}

\begin{question}

Betekent dit dat er geen gravitatiekrachten op $B$ werken?

\end{question}

\begin{question}

Geef voor $\vec F_{A\rightarrow B}$ het bijhorende actie-reactiepaar.

\end{question}

\begin{question}

Geef voor $\vec F_{C\rightarrow B}$ het bijhorende actie-reactiepaar.

\end{question}


\begin{oplossing}

\begin{enumerate}

  \item De kracht van A op B wijst naar links en de kracht van C op B

  naar rechts. Omdat de krachtgroottes gelijk zijn:

  \[\vec F_{A\rightarrow B} + \vec F_{C\rightarrow B} = \vec0.\]

  Dus:

  \[\boxed{\vec F_{\mathrm{res},B}=\vec0}.\]

  \item Volgens Newton II:

  \[\sum\vec F=m_B\vec a_B.\]

  Omdat de resulterende kracht nul is:

  \[\boxed{\vec a_B=\vec0}.\]

  \item Nee. Er werken wel twee gravitatiekrachten op B. Ze zijn alleen

  even groot en tegengesteld gericht, zodat hun vectorsom nul is.

  \item Bij

  \[\vec F_{A\rightarrow B}\]

  hoort volgens Newton III:

  \[\boxed{ \vec F_{B\rightarrow A} = -\vec F_{A\rightarrow B} }.\]

  \item Bij

  \[\vec F_{C\rightarrow B}\]

  hoort:

  \[\boxed{ \vec F_{B\rightarrow C} = -\vec F_{C\rightarrow B} }.\]

\end{enumerate}

\end{oplossing}

\end{oefening}


\begin{oefening}[Van gravitatie naar elektriciteit]{\oefG}

Beschouw eerst twee massa's die elkaar door gravitatie aantrekken.

Beschouw daarna twee positieve elektrische ladingen die elkaar

afstoten.

Maak voor beide situaties een schets met de twee krachtvectoren.

Onderzoek vervolgens:

\begin{enumerate}

 \item wat in beide situaties hetzelfde is volgens Newton III;

 \item wat verschillend is in de zin van de krachtvectoren;

 \item waarom in beide gevallen de twee krachten even groot moeten zijn;

 \item waarom de versnellingen van beide systemen toch verschillend

    kunnen zijn.

\end{enumerate}


\begin{oplossing}

\begin{enumerate}

  \item In beide situaties geldt Newton III.

  Voor gravitatie:

  \[\vec F_{A\rightarrow B} = -\vec F_{B\rightarrow A}.\]

  Voor de elektrische interactie geldt dezelfde structuur:

  \[\vec F_{1\rightarrow2} = -\vec F_{2\rightarrow1}.\]

  De twee krachten zijn dus telkens even groot, liggen langs dezelfde

  werklijn en hebben tegengestelde zin.

  \item Bij gravitatie tussen twee massa's is de wisselwerking

  aantrekkend. De krachtvectoren wijzen dus naar elkaar toe.

  Twee positieve elektrische ladingen stoten elkaar af. De

  krachtvectoren wijzen daarom van elkaar weg.

  \item In beide gevallen gaat het om één interactie tussen twee

  systemen. Volgens Newton III moeten de twee krachten van die

  interactie even groot en tegengesteld gericht zijn.

  \item Voor elk systeem afzonderlijk geldt:

  \[a=\frac{F}{m}.\]

  De krachtgroottes kunnen dus gelijk zijn terwijl de versnellingen

  verschillend zijn wanneer de massa's verschillen.

  Voor twee systemen A en B:

  \[\frac{a_A}{a_B} = \frac{m_B}{m_A}.\]

\end{enumerate}

\end{oplossing}

\end{oefening}


\end{document}